%------------------------------------------------------------------------------%
\begin{question}
  Mostre que os planos
  \[
    \pi_{1}\colon x + y + z - 3 = 0
  \]
  \[
    \pi_{2}\colon 2x - y + z - 1 = 0
  \]
  são secantes e determine uma equação paramétrica para a reta de interseção
\end{question}

%------------------------------------------------------------------------------%
\begin{resolution}{Solução}
  \onslide<+->{Os vetores normais são
  \[
    \vec{n}_{1} = \vetortri{1}{1}{1}
    \qquad\qquad
    \vec{n}_{2} = \vetortri{2}{-1}{1}
  \]}
  \onslide<+->{Não são paralelos, pois}
  \begin{align*}
    \onslide<+->{1 & = 1} \\
    \onslide<+->{2 & \neq -1}
  \end{align*}
  \onslide<+->{Os planos são secantes}
\end{resolution}

%------------------------------------------------------------------------------%
\begin{resolution}{Solução}
  \onslide<+->{A direção da reta de interseção é}
  \begin{align*}
    \onslide<+->{\vec{v} & = \vec{n}_{1}\times\vec{n}_{2} \\}
    \onslide<+->{& = \begin{vmatrix}
      \vec{i} & \vec{j} & \vec{k} \\
      1       & 1       & 1       \\
      2       & -1      & 1
    \end{vmatrix}}
    \\
    \onslide<+->{& = (2, 1, -3)}
  \end{align*}
\end{resolution}

%------------------------------------------------------------------------------%
\begin{resolution}{Solução}
\begin{columns}[t]
\column{0.4\textwidth}
  \onslide<+->{Ponto comum aos dois planos}
  \onslide<+->{\[
    \begin{cases}
       x + y + z = 3  \\
      2x - y + z = 1
    \end{cases}
  \]}
  \onslide<+->{Escolhemos $z = 0$}
  \onslide<+->{\[
    \begin{cases}
       x + y = 3  \\
      2x - y = 1
    \end{cases}
  \]}
\column{0.3\textwidth}
  \onslide<+->{Somando as equações}
  \begin{align*}
    \onslide<+->{3x & = 4 \\}
    \onslide<+->{x  & = \frac{4}{3} \\}
    \onslide<+->{y  & = 3-x \\}
    \onslide<+->{   & = 3 - \frac{4}{3}\\}
    \onslide<+->{   & = \frac{5}{3}}
  \end{align*}
\column{0.3\textwidth}
  \onslide<+->{Um ponto da reta é
  \[
    A = \left( \frac{4}{3}, \frac{5}{3}, 0 \right)
  \]}
\end{columns}
\end{resolution}

%------------------------------------------------------------------------------%
\begin{resolution}{Solução}
\begin{columns}[t]
\column{0.5\textwidth}
  \onslide<+->{Equação vetorial da reta de interseção}
  \onslide<+->{\[
    \vetortri{x}{y}{z}
    = \Vetortri{\frac{4}{3}}{\frac{5}{3}}{0}
    + t \vetortri{2}{1}{-3}
  \]}
\column{0.5\textwidth}
  \onslide<+->{A equação paramétrica é}
  \onslide<+->{
  \[
    \begin{cases}
      x = \dfrac{4}{3} + 2t \\[5mm]
      y = \dfrac{5}{3} +  t \\[5mm]
      z = -3t
    \end{cases}
  \]}
\end{columns}
\end{resolution}

%------------------------------------------------------------------------------%
\endinput


